\documentclass{article}
\newtheorem{theorem}{Theorem}
\begin{document}
A family has tail fade when, for every tolerance $\varepsilon>0$, there is an index $N$ such that
\[
 |u_n(x)|<\varepsilon \qquad n\geq N.
\]
A fixed horizon $H>0$ governs each comparison.
An historical aside discusses an unrelated normalization convention.
\begin{theorem}\label{thm:tail}
Every family with tail fade satisfies the uniform estimate on the fixed horizon.
\end{theorem}
\begin{proof}
Use tail fade beyond its witness index and keep the fixed horizon throughout the estimate.
\end{proof}
\end{document}
